\begin{resumo}[Abstract]
 \begin{otherlanguage*}{english}
Outsourcing software development in the Brazilian Public Administration is widely adopted as a strategy to optimize costs and transfer the operational burden to the private sector. However, delegating development does not exempt the public administration from responsibility for the quality of the delivered product, whose deficiencies directly impact citizens. The current regulatory framework, comprising Law No. 14,133/2021 and Ordinance SGD/MGI No. 750/2023, establishes guidelines for these procurements, but delegates the definition of quality assessment models, metrics, and criteria to the contract manager—resulting in a critical operational gap for the technical inspection of deliverables. In this context, this work empirically investigates how to operationalize Maintainability quality attributes, specifically Modularity and Testability, through objective indicators capable of supporting the acceptance of software product releases developed by third parties. To achieve this, a Structured Literature Review was conducted to gather bibliographic groundings, alongside the planning of a case study for quantitative evaluation. This evaluation is performed through the automated extraction of metrics from the history of repository-integrated releases, aggregated according to the principles of the MeasureSoftGram meta-model.

\vspace{\onelineskip}

\noindent
\textbf{Keywords}: software quality. maintainability. outsourcing. public contracts. code metrics. static analysis. case study.
 \end{otherlanguage*}
\end{resumo}
